% udt — example of the script mode
\documentclass[a4paper, polish, 11pt]{scrartcl}
\usepackage[script]{udt}   % [script, full] — all variants

\newvariant{supplement}   % conditional content, see below
\showvariant{supplement}  % comment out for the short version

\institution{Politechnika Gdańska, Wydział FTiMS}
\course{Mechanika klasyczna}
\group{Fizyka Techniczna, III rok}
\title{Moment bezwładności}
\subtitle{Przykład dokumentu w trybie script}
\author{Jan Kowalski}
\studentid{000000}
\place{Gdańsk}
\date{30 września 2026}

\begin{document}
\maketitle

\begin{abstract}
Przykład pokazuje blok tytułowy, spis treści, środowiska twierdzeń
i odsyłacze.
\end{abstract}

\tableofcontents

\section{Definicja}\label{sec:def}
\begin{definition}[Moment bezwładności]\label{def:I}
Momentem bezwładności bryły względem osi o wersorze $\vu{n}$ nazywamy
\begin{equation}
	I = \int h^2 \dd m, \qquad h^2 = \rho^2 - (\vu{n} \vdot \vb{\rho})^2
	\label{eq:I}
\end{equation}
\end{definition}

\section{Ślad tensora bezwładności}\label{sec:slad}
\begin{theorem}\label{tw:slad}
Suma momentów względem trzech wzajemnie prostopadłych osi
przecinających się w punkcie $O$ wynosi
$I_x + I_y + I_z = 2 \int r^2 \dd m$.
\end{theorem}
\begin{proof}
Dla osi $x$ mamy $h^2 = y^2 + z^2$; sumując po trzech osiach, każda
współrzędna pojawia się dwukrotnie.
\end{proof}

\begin{fact}\label{f:niezm}
Prawa strona tej równości nie zależy od orientacji osi.
\end{fact}

\begin{supplement}
\begin{remark}
Ta uwaga jest w wariancie \texttt{supplement}: znika, gdy usunie się
\verb|\showvariant{supplement}|.
\end{remark}
\end{supplement}

\begin{extended}
\begin{aside}[Wariant extended]
Blok widoczny tylko z opcją \texttt{full} albo po
\verb|\showvariant{extended}|.
\end{aside}
\end{extended}

Zob. \cref{def:I}, \cref{eq:I}, \cref{tw:slad} oraz \cref{sec:def,sec:slad}.

\end{document}
